A salesperson should always have three consistent goals\footcite{care_2014_mastering}:
\begin{enumerate}
    \item Become obsessed with the process of becoming better every single day in work (doing something significant).
    \item Care for your family, and friends (in love).
    \item Be courageous even in the most difficult of times.
\end{enumerate}

You'll notice, my interpretation is different from Care's in \textit{Mastering Technical Sales}. He notes that making money should be a goal, but I disagree with him. Money isn't everything, nor should it be the goal. I recall reading \textit{Man's Search for Meaning} in my early 20s, and coming across the Preface, which was written in 1992 by Viktor E. Frankl while he was staying in Vienna.

In a passage, he wrote:

\begin{quote}
    ``Don't aim at success---the more you aim at it and make it a target, the more you are going to miss it. For success, like happiness, cannot be pursued; it must ensue, and it only does so as the unintended side-effect of one's dedication to a cause greater than oneself or as the by-product of one's surrender to a person other than oneself. Happiness must happen, and the same holds for success: you have to let it happen by not caring about it. I want you to listen to what your conscience commands you to do and go on to carry it out to the best of your knowledge. Then you will live to see that in the long run---in the long run, I say!---success will follow you precisely because you had forgotten to think of it.''\footcite{frankl_2017_mans}
\end{quote}

Now, I want to break this passage's main messaging down:

\begin{itemize}
    \item Don't chase success directly.
    \begin{itemize}
        \item If you end up making success with all you care about, you will miss it. Success must ensue with a process that is ridden with delusional self-belief.
    \end{itemize}

    \item Success and happiness are side effects.
    \begin{itemize}
        \item Success and happiness come when you're busy doing something you find meaningful.
    \end{itemize}

    \item Listen to your conscience.
    \begin{itemize}
        \item Do what you believe is right, not what will make you successful or happy.
    \end{itemize}

    \item Focus on something bigger than you.
    \begin{itemize}
        \item Instead of trying to become something, find a purposed goal that is bigger than you.
        \begin{itemize}
            \item For example, an ideal cancer researcher does not go to work worrying about how their work will make them happy, or how it will benefit their personal life. Instead, they have a pursuit of solving the problem at the end, which is limiting, and maybe one day, finding a cure for cancer.
        \end{itemize}
    \end{itemize}

    \item Let go of the outcome.
    \begin{itemize}
        \item If you are on the right track, you need to find a purposed approach to your pursuits that is void of being obsessed with the outcome.
        \item For instance, suppose I was aiming to become a \textit{bigshot} entrepreneur. I more likely than not will fail multiple times in my life. I might never become a successful entrepreneur. However, being able to take solace in the fact that outside of that outcome, you can still maintain some peace of semblance.
        \item Another example could be death. We are all going to die. So, what is the point? The point is not to give up because we are going to die, the point is to try even harder because a life lived unfulfilled is a life not lived at all.
    \end{itemize}
\end{itemize}